\textbf{4.\ (Cake division).}

Albert and Bernie divide a cake by alternating proposals. First Albert proposes
a division. If Bernie accepts, the division is implemented. Otherwise they wait
one hour and Bernie proposes. If Albert accepts, the division is implemented;
otherwise they wait another hour, and so on.

A fraction \(p\) of the cake is worth
\[
p\left(\frac12\right)^t
\]
after \(t\) hours of waiting, for \(t\le 5\). After \(5\) hours the cake becomes
worthless. Find a subgame perfect Nash equilibrium and determine how long they wait.

\textbf{Solution.}

The last possible valuable proposal is at \(t=5\), made by Bernie. If Albert
rejects at \(t=5\), both get \(0\). Hence Bernie can keep the whole remaining cake:
\[
V_5^P=\left(\frac12\right)^5=\frac1{32}
\]

Applying backward induction. At time \(t\), the responder can reject and become proposer at
time \(t+1\). Therefore the responder must receive exactly the value with which the responder would accept the proposal
\[
V_{t+1}^P -
\]

Hence
\[
V_t^P=\left(\frac1{2}\right)^t - V_{t+1}^P
\]

Now compute with \(V_6=0\):
\[
V_5^P=\left(\frac1{2}\right)^5 - V_6=\frac1{32} - 0=\frac1{32},
\]
\[
V_4^P=\left(\frac1{2}\right)^4 - V_5=\frac1{16} - \frac1{32}=\frac1{32},
\]
\[
V_3^P=\left(\frac1{2}\right)^3 - V_4=\frac1{8} - \frac1{16}=\frac3{32},
\]
\[
V_2^P=\left(\frac1{2}\right)^2 - V_3=\frac1{4} - \frac3{32}=\frac5{32},
\]
\[
V_1^P=\left(\frac1{2}\right)^1 - V_2=\frac1{2} - \frac5{32}=\frac{11}{32},
\]
\[
V_0^P=\left(\frac1{2}\right)^0 - V_1=\frac{21}{32}
\]

At \(t=0\), Albert is proposer. He gives Bernie her continuation value:
\[
\frac{11}{32}
\]

Therefore Albert proposes
\[
\left(\frac{21}{32},\frac{11}{32}\right),
\]
and Bernie accepts.

Thus the subgame perfect Nash equilibrium outcome is
\[
\left(\frac{21}{32},\frac{11}{32}\right)
\]
and they wait
\[
0\text{ hours}
\]
